%% SYNTHETIC TEST UNIT generated by writing/tools/make_outline_book.py from FINAL_OUTLINE_v2 -- not book text
\chapter{DIS 的全阶动量空间证明、重夸克与核靶}\label{ch:disall}
本单元为测试用的合成文本（由 FINAL\_OUTLINE\_v2 自动生成），陈述内容从略。

\section{第 1 节（测试）}\label{sec:disall-regions}

\begin{definition}[Breit 系与投影算符]\label{def:disall-breit}
（测试用合成陈述：内容从略。）
\uses{def:dis-kinematics}
\end{definition}

\begin{proposition}[从非物理区前向振幅到矩与 $x$ 空间的两条途径 \tagR]\label{prop:disall-two-routes}
（$R_{\mathrm{cond}}$：在所列编号假设之下证明完整；结论的地位不高于其中最弱的假设。）
前提：\cref{asm:appN-collins}、\cref{asm:appN-contour}、\cref{asm:appN-envelope}、\cref{asm:appN-separation}、\cref{asm:appN-cutsum}。
（测试用合成陈述：内容从略。）
\uses{def:disall-breit, def:dis-compton, def:method-collins-conditions, asm:appN-collins, asm:appN-cutsum, asm:appN-contour, asm:appN-envelope, asm:appN-separation, thm:appN-subtraction, lem:appN-error-table, lem:method-local-accuracy, lem:method-meets}
\end{proposition}

\begin{lemma}[末态喷注的非钳夹性 \tagA]\label{lem:disall-fsi}
（测试用合成陈述：内容从略。）
\uses{prop:method-fsi, def:disall-breit, prop:disall-two-routes}
\end{lemma}
\begin{proofsketch}测试。\end{proofsketch}
\paragraph{地位与缺口}测试：本陈述的地位由\cref{asm:appN-cutsum}决定。

\section{第 2 节（测试）}\label{sec:disall-fsi}

\begin{proposition}[DIS 的领头区域 \tagA]\label{prop:disall-regions}
（测试用合成陈述：内容从略。）
\uses{thm:ir-leading-regions, lem:disall-fsi}
\end{proposition}
\begin{proofsketch}测试。\end{proofsketch}
\paragraph{地位与缺口}测试：本陈述的地位由\cref{asm:appN-envelope}、\cref{asm:appN-separation}决定。

\begin{proposition}[DIS 减除项的因子化结构 \tagR]\label{prop:disall-factorized-subtraction}
（$R_{\mathrm{cond}}$：在所列编号假设之下证明完整；结论的地位不高于其中最弱的假设。）
前提：\cref{asm:appN-ward-product}、\cref{asm:appN-uv}、\cref{prop:disall-regions}。
（测试用合成陈述：内容从略。）
\uses{prop:method-product, def:method-subtraction, prop:disall-regions, def:pdf-quark, def:pdf-gluon, asm:appN-ward-product, asm:appN-uv, lem:appN-algebra}
\end{proposition}

\begin{proposition}[(H2) 的验证 \tagA]\label{prop:disall-h2}
（测试用合成陈述：内容从略。）
\uses{def:method-collinear-approx, lem:method-ward-multi, prop:method-collinear-wilson, prop:method-ward-tree, prop:disall-regions}
\end{proposition}
\begin{proofsketch}测试。\end{proofsketch}
\paragraph{地位与缺口}测试：本陈述的地位由\cref{asm:appN-ward-product}决定。

\section{第 3 节（测试）}\label{sec:disall-subtraction}

\begin{proposition}[(H3)(H4) 平凡成立 \tagA]\label{prop:disall-h3h4}
（测试用合成陈述：内容从略。）
\uses{prop:disall-regions, lem:disall-fsi, lem:method-soft-unitarity}
\end{proposition}
\begin{proofsketch}测试。\end{proofsketch}
\paragraph{地位与缺口}测试：本陈述的地位由\cref{lem:disall-fsi}、\cref{prop:disall-regions}决定。

\begin{theorem}[减除与 $\overline{\rm MS}$ 重整化的相容性（矩空间） \tagR]\label{thm:disall-renorm-compat}
（$R_{\mathrm{cond}}$：在所列编号假设之下证明完整；结论的地位不高于其中最弱的假设。）
前提：\cref{asm:appN-multren}、\cref{asm:appN-ir-regulator}、\cref{asm:appN-contour}、\cref{prop:disall-two-routes}、\cref{prop:disall-factorized-subtraction}。
（测试用合成陈述：内容从略。）
\uses{prop:disall-factorized-subtraction, lem:method-no-pole, prop:pdf-partonic-dr, def:prelim-msbar, asm:appN-multren, asm:appN-ir-regulator, asm:appN-contour, prop:disall-two-routes, cor:method-removing-eps}
\end{theorem}

\begin{proposition}[(H5) 的验证 \tagA]\label{prop:disall-h5}
（测试用合成陈述：内容从略。）
\uses{prop:pdf-partonic-dr, thm:method-subtraction, prop:disall-factorized-subtraction, asm:appN-multren, asm:appN-collins}
\end{proposition}
\begin{proofsketch}测试。\end{proofsketch}
\paragraph{地位与缺口}测试：本陈述的地位由\cref{asm:appN-collins}、\cref{asm:appN-multren}、\cref{prop:disall-factorized-subtraction}、\cref{thm:disall-renorm-compat}、\cref{asm:appN-ir-regulator}决定。

\section{第 4 节（测试）}\label{sec:disall-theorem}

\begin{example}[两级梯形图：三个区域与分组余项]\label{ex:disall-ladder}
（测试用合成陈述：内容从略。）
\uses{thm:method-subtraction, prop:method-chains, def:method-collins-conditions, thm:disall-renorm-compat, prop:appN-ladder}
\end{example}

\begin{theorem}[DIS 共线因子化（动量空间，$x$ 空间逐点） \tagA]\label{thm:disall-factorization}
（测试用合成陈述：内容从略。）
非微扰地位：\cref{asm:intro-bridge}。
\uses{thm:method-template, prop:disall-regions, lem:disall-fsi, prop:disall-h2, prop:disall-h3h4, prop:disall-h5, prop:disall-two-routes, def:pdf-quark, def:pdf-gluon, asm:intro-bridge, asm:appN-collins, asm:appN-envelope, asm:appN-separation, asm:appN-cutsum, asm:appN-ward-product, asm:appN-uv, asm:appN-multren}
\end{theorem}
\begin{proofsketch}测试。\end{proofsketch}
\paragraph{地位与缺口}测试：本陈述的地位由\cref{asm:appN-collins}、\cref{asm:appN-envelope}、\cref{asm:appN-separation}、\cref{asm:appN-cutsum}、\cref{asm:appN-ward-product}、\cref{asm:appN-uv}、\cref{asm:appN-multren}、\cref{asm:intro-bridge}、\cref{asm:appN-brst-nonlocal}、\cref{lem:method-ward-multi}、\cref{thm:method-template}决定。

\begin{corollary}[与 OPE 系数的一致性 \tagR]\label{cor:disall-consistency}
（$R_{\mathrm{OPE}}$：以 OPE 为前提而严格。QCD 中含夸克的微扰 OPE 是全阶论证（\cref{thm:dis-ope}），严格证明只对标量理论与纯 Yang--Mills 的规范不变算符给出；对强子态还须非微扰 OPE（\cref{asm:intro-np-ope}）。）
（测试用合成陈述：内容从略。）
\uses{thm:pdf-dis-xspace, prop:dis-wilson-partons, cor:appK-pdf, prop:pdf-partonic-dr}
\end{corollary}

\section{第 5 节（测试）}\label{sec:disall-consistency}

\begin{proposition}[$O(\alpha_s)$ 的 $\overline{\rm MS}$ 系数函数 \tagR]\label{prop:disall-nlo}
（$R_{\mathrm{PT}}$：在重整化微扰论中对红外正规化的部分子态严格。）
（测试用合成陈述：内容从略。）
（本命题只涉及部分子态。）
\uses{def:disall-breit, prop:sudakov-oneloop, lem:prelim-plus-expansion, prop:pdf-partonic-dr, lem:appC-phase-space, prop:dis-c-oneloop}
\end{proposition}

\begin{example}[$O(\alpha_s)$ 系数函数的完整推导]\label{ex:disall-nlo}
（测试用合成陈述：内容从略。）
\uses{prop:disall-nlo}
\end{example}

\begin{example}[光子—胶子融合]\label{ex:disall-gluon}
（测试用合成陈述：内容从略。）
\uses{prop:disall-nlo}
\end{example}

\section{第 6 节（测试）}\label{sec:disall-nlo}

\begin{proposition}[标度无关性 \tagR]\label{prop:disall-scale}
（$R_{\mathrm{PT}}$：在重整化微扰论中对红外正规化的部分子态严格。）
（测试用合成陈述：内容从略。）
\uses{prop:disall-nlo, thm:pdf-dglap, prop:prelim-beta}
\end{proposition}

\begin{proposition}[DIS 方案 \tagR]\label{prop:disall-dis-scheme}
（$R_{\mathrm{PT}}$：在重整化微扰论中对红外正规化的部分子态严格。）
（测试用合成陈述：内容从略。）
\uses{prop:disall-nlo, rem:pdf-schemes}
\end{proposition}

\begin{theorem}[含重夸克质量的因子化（Collins 1998） \tagA]\label{thm:disall-heavy}
（测试用合成陈述：内容从略。）
非微扰地位：\cref{asm:intro-bridge}。
\uses{thm:disall-factorization, rem:qcd-decoupling, asm:intro-bridge}
\end{theorem}
\begin{proofsketch}测试。\end{proofsketch}
\paragraph{地位与缺口}测试：本陈述的地位由\cref{thm:disall-factorization}决定。

\section{第 7 节（测试）}\label{sec:disall-scale}

\begin{example}[领头阶 $\gamma^*g\to Q\bar Q$]\label{ex:disall-heavy-lo}
（测试用合成陈述：内容从略。）
\uses{thm:disall-heavy}
\end{example}

\begin{proposition}[核靶上的 DIS \tagA]\label{prop:disall-nuclear}
（测试用合成陈述：内容从略。）
非微扰地位：\cref{asm:intro-bridge}。
\uses{thm:disall-factorization, asm:intro-bridge}
\end{proposition}
\begin{proofsketch}测试。\end{proofsketch}
\paragraph{地位与缺口}测试：见文献。
